\documentclass[11pt,a4paper]{article}

\usepackage[margin=2.2cm]{geometry}
\usepackage{microtype}
\usepackage{enumitem}
\usepackage{parskip}
\usepackage[hidelinks]{hyperref}
\usepackage{titlesec}

\titleformat{\section}{\Large\bfseries}{}{0pt}{}
\titleformat{\subsection}{\large\bfseries}{}{0pt}{}

\setlist[enumerate]{leftmargin=2.2em,itemsep=0.15em}
\setlist[itemize]{leftmargin=1.6em,itemsep=0.15em}

\title{Criterion A: Problem Specification}
\author{}
\date{}

\begin{document}

\section*{Criterion A: Problem Specification}

\subsection*{Problem scenario}

As an IB Diploma Programme Computer Science student, I and my classmates regularly write, revise, and review pseudocode in the IB-specific notation (\texttt{METHOD}, \texttt{loop while}, \texttt{output}, \texttt{input}). Despite pseudocode being central to the syllabus and to Paper~1 and Paper~2 assessment, no online editor targets the IB dialect: Replit, CodePen, and web-based IDEs all assume Python, JavaScript, or another mainstream language, and provide no parser or interpreter for IB notation. The current peer-review workflow is to paste plain text into Google~Docs or Discord, where there is no execution, no syntax checking, and no concept of simultaneous shared editing. This problem relates directly to the world around me --- collaborative learning in a current educational setting where digital tooling has not caught up with the curriculum --- and to a wider issue in computing, where mature real-time collaborative editing (Google~Docs, VS~Code Live~Share) and per-document version control (Git, Notion) exist for general-purpose work but are absent from education-focused pseudocode environments. The proposed solution is a web-based platform on which students create accounts, organise pseudocode files into workspaces, share their work publicly or with named collaborators, edit the same file simultaneously, and trace per-file revision history.

\subsection*{Computational context}

The platform is a full-stack web application written in \textbf{TypeScript}, chosen for static typing across both client and server so that schema, API contracts, and UI props share a single source of truth. \textbf{Next.js~14} (App Router) provides unified server-side rendering and React client components, removing the need for a separately deployed back end. Persistence uses \textbf{PostgreSQL} on Neon's serverless tier through the \textbf{Prisma} ORM, which generates a fully typed query client from one schema. Real-time co-editing uses \textbf{Yjs} CRDTs over WebSockets, since CRDTs guarantee convergence without server-side conflict resolution; the WebSocket server runs on \textbf{Fly.io} because Vercel (host of the Next.js app) cannot maintain long-lived sockets. Authentication uses bcrypt password hashing and signed JWTs in \texttt{httpOnly} cookies. The IB-pseudocode interpreter is implemented client-side in JavaScript to avoid network round-trips during program execution.

\begin{center}
\begin{tabular}{|p{3.4cm}|p{3.0cm}|p{4.0cm}|p{5.4cm}|}
\hline
\textbf{Context} & \textbf{Chosen} & \textbf{Alternatives considered} & \textbf{Justification} \\
\hline
Main language & TypeScript & Python (Django), Java (Spring) & Type-safety across client/server; shared models; large ecosystem for web. \\
\hline
Framework & Next.js 14 & Express + React SPA, SvelteKit & Unified SSR + client components; serverless-friendly; single deployment unit. \\
\hline
Database / ORM & PostgreSQL + Prisma & MongoDB, SQLite + raw SQL & Relational fit (users, workspaces, files, versions); typed queries; free Neon tier. \\
\hline
Realtime engine & Yjs + WebSocket & Operational Transform, polling, Firebase Realtime DB & CRDT guarantees convergence; self-hostable; no vendor lock-in. \\
\hline
Hosting & Vercel + Fly.io & Single VPS (AWS / DigitalOcean) & Vercel free for Next.js; Fly.io hosts long-lived WebSockets that Vercel cannot. \\
\hline
Data environment & Neon Postgres & Supabase, local Postgres & Free serverless tier; auto-scaling; SQL standard. \\
\hline
\end{tabular}
\end{center}

\subsection*{Success criteria}

The product will be considered successful if it meets the following:

\begin{enumerate}
  \item Users can register with an email address and password, receive a verification email, and log in only after the address is verified.
  \item The system parses and executes valid IB-pseudocode files, displaying program output and runtime errors with the offending line number.
  \item Each user can create, rename, and delete multiple workspaces, each containing multiple pseudocode files with the \texttt{.psc} extension.
  \item Workspaces can be marked public (any visitor with the URL can view) or private (only the owner and explicitly invited collaborators can view and edit).
  \item Two or more collaborators editing the same file see each other's edits propagated within 500\,ms under typical typing speeds, with no lost or duplicated characters.
  \item Every file maintains an automatic per-version history; users can view any previous version with a visual diff against the current state and restore it.
  \item Any logged-in user can fork a public workspace, producing an independent editable copy under their own account.
  \item The interface is responsive and remains usable on a mobile browser at viewport widths $\geq$~360\,px.
  \item Passwords are stored hashed (bcrypt, cost factor $\geq 10$) and session tokens are kept in \texttt{httpOnly} cookies inaccessible to client JavaScript.
  \item The main dashboard loads in under 3 seconds on a typical broadband connection ($\geq$~10\,Mbps).
\end{enumerate}

\vspace{0.4em}
\noindent\textit{Word count (problem scenario + computational context): 309.}

\end{document}
